
\documentclass[final]{elsarticle}

\usepackage[margin=2.50cm]{geometry}  % tweak margins
%\usepackage{draftcopy}                % uncomment to print a "DRAFT-COPY" in background
\usepackage{graphicx}                 % if you want to include PostScript figures
\usepackage[figuresright]{rotating}   % if you have landscape tables
\usepackage{multirow}                 % fancy tables
\usepackage[toc,page]{appendix}       % appendix
\usepackage{lineno}                   % for printing the line numbers
\usepackage{natbib}                   % for building bibliography with \cite{}
\usepackage{setspace}
\usepackage{times}
\usepackage{afterpage}
\usepackage{float}
\usepackage{subcaption}
\usepackage{longtable}
\usepackage{blkarray}                 % for correlation matrices
\usepackage{color}
\usepackage{rotate}                   % rotate text in tables
\usepackage{lscape}                   % allow page to be landscape for large tables
\usepackage{amsfonts}                 % extra symbols
\usepackage{amsmath}
%\usepackage{floatpag}                 % \thisfloatpagestyle{empty}, to suppress page numbers
\usepackage{lmodern}
\usepackage{dcolumn}% Align table columns on decimal point
\usepackage{caption}% Use captionof to be able to minipage figure's \& tables together, but still get correct captions
\usepackage{enumitem}
\usepackage[breaklinks=true,hyperfigures=true]{hyperref} % make contents and references links, but do not colour them

\setenumerate{listparindent=\parindent} % Use paragraph indentation within the enumerate
\usepackage{cleveref}
% use to fix date
% \def \today { February 1st, 2009 }
% put your own definitions here:
\newcommand{\ck}{$\checkmark$}
\newcommand{\ttbs}{\char'134}
%\newcommand{\AmS}{{\protect\the\textfont2
%  A\kern-.1667em\lower.5ex\hbox{M}\kern-.125emS}}
% add words to TeX's hyphenation exception list
%\hyphenation{author another created financial paper re-commend-ed Post-Script}

\usepackage{listings}   
\lstset{language=C++}

% Particle aliases
\newcommand{\nue}{\nu_{e}}
\newcommand{\numu}{\nu_{\mu}}
\newcommand{\nutau}{\nu_{\tau}}
\newcommand{\nueb}{\xbar{\nu}_{e}}
\newcommand{\numub}{\xbar{\nu}_{\mu}}
\newcommand{\nutaub}{\xbar{\nu}_{\tau}}
\newcommand{\nub}{\xbar{\nu}}
\newcommand{\nua}{\nu_{\alpha}}
\newcommand{\nuab}{\xbar{\nu_{\alpha}}}

\newcommand{\numuany}{\overset{(-)}{\nu_\mu}}
\newcommand{\nueany}{\overset{(-)}{\nu_e}}


% By default, LaTeX doesn't like to fill more than 0.7 of a text page with tables and graphics, 
% nor does it like too many figures per page. This behaviour can be changed:
\renewcommand\floatpagefraction{.9}
\renewcommand\topfraction{.9}
\renewcommand\bottomfraction{.9}
\renewcommand\textfraction{.1}   
\setcounter{totalnumber}{50}
\setcounter{topnumber}{50}
\setcounter{bottomnumber}{50}

%define new fonts, smaller than tiny (for showing large covariance matrices in a page)
\makeatletter
\ifcase \@ptsize \relax% 10pt
  \newcommand{\miniscule}{\@setfontsize\miniscule{4}{5}}% \tiny: 5/6
  \newcommand{\verytiny}{\@setfontsize\miniscule{3}{4}}% \tiny: 5/6
\or% 11pt
  \newcommand{\miniscule}{\@setfontsize\miniscule{5}{6}}% \tiny: 6/7
  \newcommand{\verytiny}{\@setfontsize\miniscule{4}{5}}% \tiny: 6/7
\or% 12pt
  \newcommand{\miniscule}{\@setfontsize\miniscule{5}{6}}% \tiny: 6/7
  \newcommand{\verytiny}{\@setfontsize\miniscule{4}{5}}% \tiny: 6/7
\fi
\makeatother




\linenumbers

\begin{document}

Dear authors,

here below some comments from my first reading.  I might add something
going through the new version containting also F\&C results. In order
not to delay the process here are my first reactions. There is nothing
really major but rather some suggestions to improve readibility and typos.

Andrea.

$=================================================================================================$
\newline

version: http://www.t2k.org/docs/technotes/266/joint-fit-analysis-of-the-nu-and-nubar-sk-samples
\newline

31 Could you use less digits for the number of expected events (at least for the final version)?
This recommendation also applies to several tables (i.e Tab.1). I guess you are generating them automatically.

\textcolor{blue}{The numbers in the abstract have been reduced to 2 digits. In table 1 the numbers have no errors and are exactly 
the numbers used in the MC, so I would keep all the digits. The same is for the event rates, because they are also useful 
for comparison between the oscillation analyses.}

32 Super K $\rightarrow$ Super-K

\textcolor{blue}{DONE.}

32 Since the abstract should be general I would report directly the values corresponding to Asimov-A in
place of referring to Tab. 12 (maybe just $\delta_{CP}$=... and $\theta_{23}$=...)

\textcolor{blue}{DONE.}

41 -ized (use american english everywhere consistently)

\textcolor{blue}{DONE.}

41 [55] should become [11] (reorder references)

\textcolor{blue}{That citation is ordered in that way in order to keep all the reference about statistic together.
The citation has been changed with a reference to the section about statistical methods.}

89-90 is it feasible to recall shortly the pros and cons of marginalizing vs profiling to
put the note in a more general context?

\textcolor{blue}{DONE.}

94 drop "recent"

\textcolor{blue}{DONE.}

97 I would quote "BANFF" explicitly for clarity.

\textcolor{blue}{DONE.}

102 increases

\textcolor{blue}{DONE.}

107 add "." at the end.

\textcolor{blue}{DONE.}

115-116 it seems that the marginalization is peculiar of Run 1-6.

\textcolor{blue}{DONE.}

116 remove space before [55]

\textcolor{blue}{DONE.}

117 single-ring

\textcolor{blue}{DONE.}

118 add "." at the end

\textcolor{blue}{DONE.}

Sometimes I have the impression you use too many (short) paragraphs.

133 $t$

\textcolor{blue}{DONE.}

136 the use of "nominal" for the templates even if they are corrected is a bit misleading

\textcolor{blue}{DONE.}

137 the reaction mode ... $\theta$ and $\rightarrow$ repetition, I would only leave $\vec{f}$ which is introduced for the first time

\textcolor{blue}{I think it is important to repeat that the systematic parameters depend on the reconstructed variables
as well as flavor and reaction mode. I modified the sentence a little bit.}

144 of the nominal

\textcolor{blue}{DONE.}

159 "with minor tunings" sounds a bit vague

\textcolor{blue}{DONE.}

Use Tab. Ref. Sec. Fig. consistently.

\textcolor{blue}{Using cref everywhere now. DONE.}

174, 175 $I$, $L$

\textcolor{blue}{DONE.}

210 granularity $\rightarrow$ number ?

\textcolor{blue}{DONE.}

section 2.2.4 cite a reference technical note ?

\textcolor{blue}{DONE.}

219-258 the same structure is replicated for every incoming neutrino flavor, maybe this long 
list could be compactified using a placeholder for the flavour x = nu, e, ...

\textcolor{blue}{The full list of the analysis templates is useful for cross checks between the oscillation analyses
and for check with the future ones. We would like to keep the full list.}

Tab. 2 uses the same order and notation used for templates as in the list in the text

\textcolor{blue}{The list in the text has been adapted to the one in the tables.}

272-275 this is repeated, probably it is not necessary to say it again.

\textcolor{blue}{DONE.}

293 $e$-like

\textcolor{blue}{DONE.}

294 25 50-MeV bins (for consistency)

\textcolor{blue}{DONE.}

Vs $\rightarrow$ vs ? (taste)

\textcolor{blue}{Using "vs" everywhere. DONE.}

310 add (see Fig. 32 and 33) to help the reader

\textcolor{blue}{DONE.}

365, 377, 391 cross section

\textcolor{blue}{DONE.}

371, 411 cite also the new BANFF TN230 and 285 (currently [28] = TN220). It could also be useful to
explicit the TN number when quoting.

\textcolor{blue}{DONE.}

394 normalize

\textcolor{blue}{Using normalisation everywhere now. DONE.}

405 say (appendix D)

\textcolor{blue}{DONE.}

406 The fact that certain nuisance parameters affect certain templates
or not is codified in Tabs. 34, 35, 36, 37.  Is there some
ambiguity/level of arbitrarity in certain particular cases? May be
one could add a short comment on this point to help the reader (tables
are useful but very "dense").

\textcolor{blue}{Maybe it is better to leave any comment about the cross sectoin systematic
parameters to the NIWG note.
For the case of the NC$1\gamma$ normalisation parameter, the new choice was driven by a new papers in literature.}

411 section $\rightarrow$ appendix or App.

\textcolor{blue}{DONE.}

There are some repetitions which I would avoid. For example the number
of POT is given in several different places.

\textcolor{blue}{DONE.}

498 it is important

\textcolor{blue}{DONE.}

538 $n_i^{exp} = ...$ is a repetition

\textcolor{blue}{DONE.}

542 recall shortly the extended maximum likelihood or use a reference

\textcolor{blue}{DONE.}

560 drop "we" before will

\textcolor{blue}{DONE.}

560 derive $\rightarrow$ use

\textcolor{blue}{DONE.}

560 add , after on or rephrase

\textcolor{blue}{DONE.}

653 this sentence on profiling on the mass hierarchy (while marginalizing on
other parameters) creates confusion here. I would defer this to the text
later which explains better what has been done (global vs independent)

\textcolor{blue}{DONE.}

654 '' ''  $\rightarrow$ `` ''

\textcolor{blue}{DONE.}

671 Bernoulli distribution $\rightarrow$ In that paper the do not say that the delta chi2
follows a Bernoulli distribution but rather the fitted parameter, if I am not mistaken.

\textcolor{blue}{DONE.}

caption of Fig. 18: it seems wrong. Where are the comparisons with other groups ?

\textcolor{blue}{DONE.}

700-703 rephrase in a more linear way separating better Fig. 19 and 20.

\textcolor{blue}{DONE.}

Section 4. This is rather general and central. Would it be reasonable to promote it
at the beginning of the note?

\textcolor{blue}{Section 4 is now cited both in the abstract and in the introduction. DONE.}

339 has the GloBES oscillation code ever been used as a testbench ?

\textcolor{blue}{Yes it was test against GloBES as well.}

In the 1D plots of Fig. 5 the e-like final state contributions are
completely invisible even zooming in the pdf. I imagine they are
very small in the mu-like sample but could you check it ? I am not
sure if it is worth keeping them in the legend after all (not
a strong opinion on this).

\textcolor{blue}{As shown in table 4, the expected contamination from oscillate $\bar{\nu}_e$ is only 0.138
events over a total of 26.594.
I think it's good to still have it in the legend in order to be consistent with the other plots
and make it clear that it is not included in any other contribution.}

413 you say tht the correction for the bias in the fit has a small
impact on the systematics error value referring quickly to appendix D.
Maybe it could be useful saying that delta $N_SK / N_SK$ for mu-like passes
from ... to ... and also for e-like (numu or antinumu mode).

414 you concentrate on the different number of expected e-like in antineutrino
mode wrt 2015. I would also mention how it changes for neutrino mode since
some of the reasons (i.e. theta13) affects both.

\textcolor{blue}{The thing is when the neutrino mode $e$-like sample
was analyzed last time the cross section model in largely different}

Mach3JointFitRun14
The discussion starting in line 414 could be put under a subsection
("Comparison with 2015 analysis foe e-like ...")

\textcolor{blue}{DONE.}

468 Super-K

\textcolor{blue}{DONE.}

Fig. 6-16. The captions are rather long and it is difficult to
understand quickly the difference between the different inputs. Could
you highlight in bold only the relevant information or make a
reference for the caption only to the first figure? It could be also
helpful to add in the same figure both the spectra and the error
envelopes.

\textcolor{blue}{DONE.}

507 it would be useful in the discussion to also mention Tab. 42
(delta=0) besides Tab. 21 (delta = pi/2). This part is interesting but
sounds a bit qualitative. Can you reference a table/TM when describing
the correlation of sigma(nue) and sigma(anti-nue) to go in this
direction ? I remember this information in the NIWG or BANFF TNs.

\textcolor{blue}{DONE.}

1001 say in the text that Tab. 38 is for e-like and 39 for mu-like nueutrino mode
while Tab. 40 ...etc (it helps).

\textcolor{blue}{Described in the text in section D. DONE.}

The choice of parameters for Asimov-A and Asimov-B is only briefly mentioned in
the caption of Tab. 12 apparently. Even of this is not so critical, I would suggest
to add a paragraph commenting on the reason for the choice ($\delta_{CP}=0$,
$\theta_{23}$ in the other octant). In particular how did you choose the $\theta_{23}$ values
for B ? from T2K Run 1-4 best fit?

\textcolor{blue}{The Asimov A is the best-fit of the Run 1-4 data Mach3 analysis.
The Asimov B was a variation of Asimov A useful to test the sensitivity to the opposite octant,
which reduce the number of expected events in the $e$-like samples as well as the 
sensitivity to $\delta_{CP}$. This is a more conservative Asimov data set.
Both Asimov A and B were chosen by the OA group as test points.
A small paragraph has been added.}

\end{document}

